\section{Personalización de gráficos en objetos}

\subsection{Paquete de configuraciones predefinidas (\texttt{.style.use()})}

Matplotlib incluye una gran variedad de estilos. Para ver cuáles tienes disponibles en tu computadora actual, usa este comando en tu terminal de Python:

\begin{pycode}{Listar Estilos}
import matplotlib.pyplot as plt

print(plt.style.available)
\end{pycode}

Los estilos mas utilizados son:
\begin{itemize}
    \item \texttt{'ggplot'}: Usa un fondo gris claro con rejillas blancas. Es muy legible para informes técnicos.
    \item \texttt{'bmh'}: (Bayesian Methods for Hackers). Ideal para gráficos estadísticos limpios y con colores sobrios.
    \item \texttt{'seaborn-v0\_8'}: Emula la estética de Seaborn. Utiliza paletas de colores que son fáciles de distinguir para personas con daltonismo.
    \item \texttt{'dark\_background'}: Perfecto para presentaciones en pantallas OLED o dashboards nocturnos.
    \item \texttt{'grayscale'}: Para cuando el documento se va a imprimir en blanco y negro; asegura que las líneas se distingan por su tono de gris.
\end{itemize}

La utilización de \texttt{plt.style.use} tiene un efecto global. Una vez ejecutada, todos los graficos que se creen heredan ese estilo.

En código, esto se aplicaria de la siguiente manera:
\begin{pycode}{Aplicar Estilo}
import matplotlib.pyplot as plt

# Aplicamos el estilo al principio del script
plt.style.use('ggplot')

# Cualquier grafico siguiente hereda el estilo
fig, ax = plt.subplots()
ax.plot([1, 2, 3], [10, 20, 15], label="Figura con estilo heredado")
ax.legend()
plt.show()
\end{pycode}


\subsubsection{Context manager (Uso temporal)}

Si queremos que solo un gráfico tenga un estilo diferente, pero que el resto mantenga el estilo estándar se utiliza \texttt{with}, \texttt{context}.

\begin{pycode}{Estilo Temporal}
# Estilo por defecto

# Estilo temporal
with plt.style.context('dark_background'):
fig, ax = plt.subplots()
ax.plot([1, 2, 3], [10, 20, 15], label="Figura con estilo heredado")

# Estilo por defecto

\end{pycode}



\subsection{Personalización de cada elemento (\texttt{plt.rcParams()})}

Este permite personalizar cada gráfico por elemento. Runtime Configuration (configuración en tiempo de ejecución) o \texttt{rc} es un objeto que contiene los valores por defecto de cada característica del gráfico.

Cada vez que se importa \textit{Matplotlib} se carga una configuración por defecto. Si se modifica el valor en \texttt{plt.rcParams}, ese cambio se vuelve el nuevo estándar para cada gráfico que se cree en esa sesión de \textit{Python}.

Para ver la cantidad inmensa de opciones que puedes modificar, puedes imprimir el objeto coplteto (son cientos):

\begin{pycode}{Ver rcParams}
import matplotlib as plt

# Ver todas las llaves disponibles para configurar
print(len(plt.rcParams.keys()))

# Buscar parámetros especificos relacionados con 'font' (fuentes)
fuentes = {k: v for k, v in plt.rcParams.items() if 'font' in k}
print(fuentes)

# Buscar parámetros especificos con 'grilla'
grid_params = {k: v for k, v in plt.rcParams.items() if 'grid' in k}
print(grid_params)
\end{pycode}

\subsubsection{Modificaciones comunes}

\emph{Tipografía y tamaños}

Hay parametros que permiten modificar el tamaño, el tipo de letra, su fuente, \dots
\begin{pycode}{Configuracion de fuentes}
plt.rcParams['font.size'] = 12          # Tamaño general
plt.rcParams['axes.titlesize'] = 16    # Tamaño del titulo del grafico
plt.rcParams['axes.labelsize'] = 14    # Tamaño de etiquetas X e Y
plt.rcParams['font.family'] = 'serif'   # Fuente con serifa (estilo Times New Roman)
\end{pycode}

\emph{Lineas y marcadores}

Para no tener que definir el grosor de cada linea, el tamaño, color, se puede definir en \texttt{plt.rcParams}.
\begin{pycode}{Configuracion de lineas}
plt.rcParams['lines.linewidth'] = 2.5   # Lineas mas gruesas
plt.rcParams['lines.markersize'] = 8     # Marcadores mas grandes
plt.rcParams['lines.color'] = 'tab:blue' # Color por defecto
\end{pycode}

\emph{El Lienzo (Figure) y la Rejilla (Grid)}

Se puede modificar el tamaño general de todos los gráficos, el estilo de grilla, su \texttt{alpha}, para que quede de forma predefinida en todos los gráficos que vaya a realizar.
\begin{pycode}{Configuracion de fondo}
plt.rcParams['figure.figsize'] = (10, 6) # Tamaño por defecto (ancho, alto)
plt.rcParams['axes.grid'] = True          # ¡Activar la grilla en todos los graficos!
plt.rcParams['grid.alpha'] = 0.3          # Hacer la grilla sutil (transparencia)
plt.rcParams['grid.linestyle'] = '--'     # Grilla punteada
\end{pycode}

\subsubsection{Forma abreviada (\texttt{plt.rc()})}

Esta forma permite cambiar muchos parámetros de un mismo \textit{grupo}. Si quisieramos cambiar todo lo referido a \texttt{fuentes}, podriamos utilizar:
\begin{verbatim}
    plt.rc('font', configuraciones\dots)
\end{verbatim}
\begin{pycode}{Metodo abreviado}
plt.rc('font', size=12, family='sans-serif', weight='bold')
plt.rc('lines', linewidth=2, linestyle='-')
\end{pycode}


\begin{pycode}{Ejeplto de personalización para este documento}
import matplotlib.pyplot as plt

# -------------------------- #
#       Personalización      #
# -------------------------- #

# ---- Para latex ---- #

plt.rcParams.update({
    "pgf.texsystem": "xelatex",    
    "font.family": "serif",
    "text.usetex": True,           
    "pgf.rcfonts": False,          
    "pgf.preamble": "\n".join([    
         r"\usepackage{fontspec}",
         r"\setmainfont{CMU Serif}",
    ])
})

# ---- Grilla ---- #

plt.rcParams.update({
    'axes.grid':True,
    'grid.linestyle':'--',
    'grid.linewidth':0.5,
    'grid.color':'gray',
    'grid.alpha':0.5
})

# ------------------ #
#       Gráfico      #
# ------------------ #

fig, ax = plt.subplots(1, 2, sharey=True, figsize=(5, 5))

# Eje 1
ax[0].plot([100, 150, 130],['Ene', 'Feb', 'Mar'] , color='blue')
ax[0].set_ylabel('Producto A')

# Eje 2
ax[1].plot([90, 110, 140], ['Ene', 'Feb', 'Mar'], color='green')
ax[1].set_ylabel('Producto B')

# Título global
fig.suptitle('Gráfico de ejes compartidos con personalización', fontsize=14)    
\end{pycode}

La visualización de personalización para este documento se ve en la figura \ref{fig:grafico_objeto_con_rcParams}
\begin{figure}[ht]
    \centering
    %% Creator: Matplotlib, PGF backend
%%
%% To include the figure in your LaTeX document, write
%%   \input{<filename>.pgf}
%%
%% Make sure the required packages are loaded in your preamble
%%   \usepackage{pgf}
%%
%% Also ensure that all the required font packages are loaded; for instance,
%% the lmodern package is sometimes necessary when using math font.
%%   \usepackage{lmodern}
%%
%% Figures using additional raster images can only be included by \input if
%% they are in the same directory as the main LaTeX file. For loading figures
%% from other directories you can use the `import` package
%%   \usepackage{import}
%%
%% and then include the figures with
%%   \import{<path to file>}{<filename>.pgf}
%%
%% Matplotlib used the following preamble
%%   \def\mathdefault#1{#1}
%%   \everymath=\expandafter{\the\everymath\displaystyle}
%%   \IfFileExists{scrextend.sty}{
%%     \usepackage[fontsize=10.000000pt]{scrextend}
%%   }{
%%     \renewcommand{\normalsize}{\fontsize{10.000000}{12.000000}\selectfont}
%%     \normalsize
%%   }
%%   \usepackage{fontspec}
%%   \setmainfont{CMU Serif}
%%   \makeatletter\@ifpackageloaded{underscore}{}{\usepackage[strings]{underscore}}\makeatother
%%
\begingroup%
\makeatletter%
\begin{pgfpicture}%
\pgfpathrectangle{\pgfpointorigin}{\pgfqpoint{4.648945in}{4.822560in}}%
\pgfusepath{use as bounding box, clip}%
\begin{pgfscope}%
\pgfsetbuttcap%
\pgfsetmiterjoin%
\definecolor{currentfill}{rgb}{1.000000,1.000000,1.000000}%
\pgfsetfillcolor{currentfill}%
\pgfsetlinewidth{0.000000pt}%
\definecolor{currentstroke}{rgb}{1.000000,1.000000,1.000000}%
\pgfsetstrokecolor{currentstroke}%
\pgfsetdash{}{0pt}%
\pgfpathmoveto{\pgfqpoint{0.000000in}{0.000000in}}%
\pgfpathlineto{\pgfqpoint{4.648945in}{0.000000in}}%
\pgfpathlineto{\pgfqpoint{4.648945in}{4.822560in}}%
\pgfpathlineto{\pgfqpoint{0.000000in}{4.822560in}}%
\pgfpathlineto{\pgfqpoint{0.000000in}{0.000000in}}%
\pgfpathclose%
\pgfusepath{fill}%
\end{pgfscope}%
\begin{pgfscope}%
\pgfsetbuttcap%
\pgfsetmiterjoin%
\definecolor{currentfill}{rgb}{1.000000,1.000000,1.000000}%
\pgfsetfillcolor{currentfill}%
\pgfsetlinewidth{0.000000pt}%
\definecolor{currentstroke}{rgb}{0.000000,0.000000,0.000000}%
\pgfsetstrokecolor{currentstroke}%
\pgfsetstrokeopacity{0.000000}%
\pgfsetdash{}{0pt}%
\pgfpathmoveto{\pgfqpoint{0.627083in}{0.320555in}}%
\pgfpathlineto{\pgfqpoint{2.321981in}{0.320555in}}%
\pgfpathlineto{\pgfqpoint{2.321981in}{4.094811in}}%
\pgfpathlineto{\pgfqpoint{0.627083in}{4.094811in}}%
\pgfpathlineto{\pgfqpoint{0.627083in}{0.320555in}}%
\pgfpathclose%
\pgfusepath{fill}%
\end{pgfscope}%
\begin{pgfscope}%
\pgfpathrectangle{\pgfqpoint{0.627083in}{0.320555in}}{\pgfqpoint{1.694898in}{3.774256in}}%
\pgfusepath{clip}%
\pgfsetbuttcap%
\pgfsetroundjoin%
\pgfsetlinewidth{0.501875pt}%
\definecolor{currentstroke}{rgb}{0.501961,0.501961,0.501961}%
\pgfsetstrokecolor{currentstroke}%
\pgfsetstrokeopacity{0.500000}%
\pgfsetdash{{1.850000pt}{0.800000pt}}{0.000000pt}%
\pgfpathmoveto{\pgfqpoint{0.704124in}{0.320555in}}%
\pgfpathlineto{\pgfqpoint{0.704124in}{4.094811in}}%
\pgfusepath{stroke}%
\end{pgfscope}%
\begin{pgfscope}%
\pgfsetbuttcap%
\pgfsetroundjoin%
\definecolor{currentfill}{rgb}{0.000000,0.000000,0.000000}%
\pgfsetfillcolor{currentfill}%
\pgfsetlinewidth{0.803000pt}%
\definecolor{currentstroke}{rgb}{0.000000,0.000000,0.000000}%
\pgfsetstrokecolor{currentstroke}%
\pgfsetdash{}{0pt}%
\pgfsys@defobject{currentmarker}{\pgfqpoint{0.000000in}{-0.048611in}}{\pgfqpoint{0.000000in}{0.000000in}}{%
\pgfpathmoveto{\pgfqpoint{0.000000in}{0.000000in}}%
\pgfpathlineto{\pgfqpoint{0.000000in}{-0.048611in}}%
\pgfusepath{stroke,fill}%
}%
\begin{pgfscope}%
\pgfsys@transformshift{0.704124in}{0.320555in}%
\pgfsys@useobject{currentmarker}{}%
\end{pgfscope}%
\end{pgfscope}%
\begin{pgfscope}%
\definecolor{textcolor}{rgb}{0.000000,0.000000,0.000000}%
\pgfsetstrokecolor{textcolor}%
\pgfsetfillcolor{textcolor}%
\pgftext[x=0.704124in,y=0.223333in,,top]{\color{textcolor}{\rmfamily\fontsize{10.000000}{12.000000}\selectfont\catcode`\^=\active\def^{\ifmmode\sp\else\^{}\fi}\catcode`\%=\active\def%{\%}$\mathdefault{100}$}}%
\end{pgfscope}%
\begin{pgfscope}%
\pgfpathrectangle{\pgfqpoint{0.627083in}{0.320555in}}{\pgfqpoint{1.694898in}{3.774256in}}%
\pgfusepath{clip}%
\pgfsetbuttcap%
\pgfsetroundjoin%
\pgfsetlinewidth{0.501875pt}%
\definecolor{currentstroke}{rgb}{0.501961,0.501961,0.501961}%
\pgfsetstrokecolor{currentstroke}%
\pgfsetstrokeopacity{0.500000}%
\pgfsetdash{{1.850000pt}{0.800000pt}}{0.000000pt}%
\pgfpathmoveto{\pgfqpoint{1.320451in}{0.320555in}}%
\pgfpathlineto{\pgfqpoint{1.320451in}{4.094811in}}%
\pgfusepath{stroke}%
\end{pgfscope}%
\begin{pgfscope}%
\pgfsetbuttcap%
\pgfsetroundjoin%
\definecolor{currentfill}{rgb}{0.000000,0.000000,0.000000}%
\pgfsetfillcolor{currentfill}%
\pgfsetlinewidth{0.803000pt}%
\definecolor{currentstroke}{rgb}{0.000000,0.000000,0.000000}%
\pgfsetstrokecolor{currentstroke}%
\pgfsetdash{}{0pt}%
\pgfsys@defobject{currentmarker}{\pgfqpoint{0.000000in}{-0.048611in}}{\pgfqpoint{0.000000in}{0.000000in}}{%
\pgfpathmoveto{\pgfqpoint{0.000000in}{0.000000in}}%
\pgfpathlineto{\pgfqpoint{0.000000in}{-0.048611in}}%
\pgfusepath{stroke,fill}%
}%
\begin{pgfscope}%
\pgfsys@transformshift{1.320451in}{0.320555in}%
\pgfsys@useobject{currentmarker}{}%
\end{pgfscope}%
\end{pgfscope}%
\begin{pgfscope}%
\definecolor{textcolor}{rgb}{0.000000,0.000000,0.000000}%
\pgfsetstrokecolor{textcolor}%
\pgfsetfillcolor{textcolor}%
\pgftext[x=1.320451in,y=0.223333in,,top]{\color{textcolor}{\rmfamily\fontsize{10.000000}{12.000000}\selectfont\catcode`\^=\active\def^{\ifmmode\sp\else\^{}\fi}\catcode`\%=\active\def%{\%}$\mathdefault{120}$}}%
\end{pgfscope}%
\begin{pgfscope}%
\pgfpathrectangle{\pgfqpoint{0.627083in}{0.320555in}}{\pgfqpoint{1.694898in}{3.774256in}}%
\pgfusepath{clip}%
\pgfsetbuttcap%
\pgfsetroundjoin%
\pgfsetlinewidth{0.501875pt}%
\definecolor{currentstroke}{rgb}{0.501961,0.501961,0.501961}%
\pgfsetstrokecolor{currentstroke}%
\pgfsetstrokeopacity{0.500000}%
\pgfsetdash{{1.850000pt}{0.800000pt}}{0.000000pt}%
\pgfpathmoveto{\pgfqpoint{1.936777in}{0.320555in}}%
\pgfpathlineto{\pgfqpoint{1.936777in}{4.094811in}}%
\pgfusepath{stroke}%
\end{pgfscope}%
\begin{pgfscope}%
\pgfsetbuttcap%
\pgfsetroundjoin%
\definecolor{currentfill}{rgb}{0.000000,0.000000,0.000000}%
\pgfsetfillcolor{currentfill}%
\pgfsetlinewidth{0.803000pt}%
\definecolor{currentstroke}{rgb}{0.000000,0.000000,0.000000}%
\pgfsetstrokecolor{currentstroke}%
\pgfsetdash{}{0pt}%
\pgfsys@defobject{currentmarker}{\pgfqpoint{0.000000in}{-0.048611in}}{\pgfqpoint{0.000000in}{0.000000in}}{%
\pgfpathmoveto{\pgfqpoint{0.000000in}{0.000000in}}%
\pgfpathlineto{\pgfqpoint{0.000000in}{-0.048611in}}%
\pgfusepath{stroke,fill}%
}%
\begin{pgfscope}%
\pgfsys@transformshift{1.936777in}{0.320555in}%
\pgfsys@useobject{currentmarker}{}%
\end{pgfscope}%
\end{pgfscope}%
\begin{pgfscope}%
\definecolor{textcolor}{rgb}{0.000000,0.000000,0.000000}%
\pgfsetstrokecolor{textcolor}%
\pgfsetfillcolor{textcolor}%
\pgftext[x=1.936777in,y=0.223333in,,top]{\color{textcolor}{\rmfamily\fontsize{10.000000}{12.000000}\selectfont\catcode`\^=\active\def^{\ifmmode\sp\else\^{}\fi}\catcode`\%=\active\def%{\%}$\mathdefault{140}$}}%
\end{pgfscope}%
\begin{pgfscope}%
\pgfpathrectangle{\pgfqpoint{0.627083in}{0.320555in}}{\pgfqpoint{1.694898in}{3.774256in}}%
\pgfusepath{clip}%
\pgfsetbuttcap%
\pgfsetroundjoin%
\pgfsetlinewidth{0.501875pt}%
\definecolor{currentstroke}{rgb}{0.501961,0.501961,0.501961}%
\pgfsetstrokecolor{currentstroke}%
\pgfsetstrokeopacity{0.500000}%
\pgfsetdash{{1.850000pt}{0.800000pt}}{0.000000pt}%
\pgfpathmoveto{\pgfqpoint{0.627083in}{0.492112in}}%
\pgfpathlineto{\pgfqpoint{2.321981in}{0.492112in}}%
\pgfusepath{stroke}%
\end{pgfscope}%
\begin{pgfscope}%
\pgfsetbuttcap%
\pgfsetroundjoin%
\definecolor{currentfill}{rgb}{0.000000,0.000000,0.000000}%
\pgfsetfillcolor{currentfill}%
\pgfsetlinewidth{0.803000pt}%
\definecolor{currentstroke}{rgb}{0.000000,0.000000,0.000000}%
\pgfsetstrokecolor{currentstroke}%
\pgfsetdash{}{0pt}%
\pgfsys@defobject{currentmarker}{\pgfqpoint{-0.048611in}{0.000000in}}{\pgfqpoint{-0.000000in}{0.000000in}}{%
\pgfpathmoveto{\pgfqpoint{-0.000000in}{0.000000in}}%
\pgfpathlineto{\pgfqpoint{-0.048611in}{0.000000in}}%
\pgfusepath{stroke,fill}%
}%
\begin{pgfscope}%
\pgfsys@transformshift{0.627083in}{0.492112in}%
\pgfsys@useobject{currentmarker}{}%
\end{pgfscope}%
\end{pgfscope}%
\begin{pgfscope}%
\definecolor{textcolor}{rgb}{0.000000,0.000000,0.000000}%
\pgfsetstrokecolor{textcolor}%
\pgfsetfillcolor{textcolor}%
\pgftext[x=0.296666in, y=0.443918in, left, base]{\color{textcolor}{\rmfamily\fontsize{10.000000}{12.000000}\selectfont\catcode`\^=\active\def^{\ifmmode\sp\else\^{}\fi}\catcode`\%=\active\def%{\%}Ene}}%
\end{pgfscope}%
\begin{pgfscope}%
\pgfpathrectangle{\pgfqpoint{0.627083in}{0.320555in}}{\pgfqpoint{1.694898in}{3.774256in}}%
\pgfusepath{clip}%
\pgfsetbuttcap%
\pgfsetroundjoin%
\pgfsetlinewidth{0.501875pt}%
\definecolor{currentstroke}{rgb}{0.501961,0.501961,0.501961}%
\pgfsetstrokecolor{currentstroke}%
\pgfsetstrokeopacity{0.500000}%
\pgfsetdash{{1.850000pt}{0.800000pt}}{0.000000pt}%
\pgfpathmoveto{\pgfqpoint{0.627083in}{2.207683in}}%
\pgfpathlineto{\pgfqpoint{2.321981in}{2.207683in}}%
\pgfusepath{stroke}%
\end{pgfscope}%
\begin{pgfscope}%
\pgfsetbuttcap%
\pgfsetroundjoin%
\definecolor{currentfill}{rgb}{0.000000,0.000000,0.000000}%
\pgfsetfillcolor{currentfill}%
\pgfsetlinewidth{0.803000pt}%
\definecolor{currentstroke}{rgb}{0.000000,0.000000,0.000000}%
\pgfsetstrokecolor{currentstroke}%
\pgfsetdash{}{0pt}%
\pgfsys@defobject{currentmarker}{\pgfqpoint{-0.048611in}{0.000000in}}{\pgfqpoint{-0.000000in}{0.000000in}}{%
\pgfpathmoveto{\pgfqpoint{-0.000000in}{0.000000in}}%
\pgfpathlineto{\pgfqpoint{-0.048611in}{0.000000in}}%
\pgfusepath{stroke,fill}%
}%
\begin{pgfscope}%
\pgfsys@transformshift{0.627083in}{2.207683in}%
\pgfsys@useobject{currentmarker}{}%
\end{pgfscope}%
\end{pgfscope}%
\begin{pgfscope}%
\definecolor{textcolor}{rgb}{0.000000,0.000000,0.000000}%
\pgfsetstrokecolor{textcolor}%
\pgfsetfillcolor{textcolor}%
\pgftext[x=0.312083in, y=2.159489in, left, base]{\color{textcolor}{\rmfamily\fontsize{10.000000}{12.000000}\selectfont\catcode`\^=\active\def^{\ifmmode\sp\else\^{}\fi}\catcode`\%=\active\def%{\%}Feb}}%
\end{pgfscope}%
\begin{pgfscope}%
\pgfpathrectangle{\pgfqpoint{0.627083in}{0.320555in}}{\pgfqpoint{1.694898in}{3.774256in}}%
\pgfusepath{clip}%
\pgfsetbuttcap%
\pgfsetroundjoin%
\pgfsetlinewidth{0.501875pt}%
\definecolor{currentstroke}{rgb}{0.501961,0.501961,0.501961}%
\pgfsetstrokecolor{currentstroke}%
\pgfsetstrokeopacity{0.500000}%
\pgfsetdash{{1.850000pt}{0.800000pt}}{0.000000pt}%
\pgfpathmoveto{\pgfqpoint{0.627083in}{3.923254in}}%
\pgfpathlineto{\pgfqpoint{2.321981in}{3.923254in}}%
\pgfusepath{stroke}%
\end{pgfscope}%
\begin{pgfscope}%
\pgfsetbuttcap%
\pgfsetroundjoin%
\definecolor{currentfill}{rgb}{0.000000,0.000000,0.000000}%
\pgfsetfillcolor{currentfill}%
\pgfsetlinewidth{0.803000pt}%
\definecolor{currentstroke}{rgb}{0.000000,0.000000,0.000000}%
\pgfsetstrokecolor{currentstroke}%
\pgfsetdash{}{0pt}%
\pgfsys@defobject{currentmarker}{\pgfqpoint{-0.048611in}{0.000000in}}{\pgfqpoint{-0.000000in}{0.000000in}}{%
\pgfpathmoveto{\pgfqpoint{-0.000000in}{0.000000in}}%
\pgfpathlineto{\pgfqpoint{-0.048611in}{0.000000in}}%
\pgfusepath{stroke,fill}%
}%
\begin{pgfscope}%
\pgfsys@transformshift{0.627083in}{3.923254in}%
\pgfsys@useobject{currentmarker}{}%
\end{pgfscope}%
\end{pgfscope}%
\begin{pgfscope}%
\definecolor{textcolor}{rgb}{0.000000,0.000000,0.000000}%
\pgfsetstrokecolor{textcolor}%
\pgfsetfillcolor{textcolor}%
\pgftext[x=0.278889in, y=3.875060in, left, base]{\color{textcolor}{\rmfamily\fontsize{10.000000}{12.000000}\selectfont\catcode`\^=\active\def^{\ifmmode\sp\else\^{}\fi}\catcode`\%=\active\def%{\%}Mar}}%
\end{pgfscope}%
\begin{pgfscope}%
\definecolor{textcolor}{rgb}{0.000000,0.000000,0.000000}%
\pgfsetstrokecolor{textcolor}%
\pgfsetfillcolor{textcolor}%
\pgftext[x=0.223333in,y=2.207683in,,bottom,rotate=90.000000]{\color{textcolor}{\rmfamily\fontsize{10.000000}{12.000000}\selectfont\catcode`\^=\active\def^{\ifmmode\sp\else\^{}\fi}\catcode`\%=\active\def%{\%}Producto A}}%
\end{pgfscope}%
\begin{pgfscope}%
\pgfpathrectangle{\pgfqpoint{0.627083in}{0.320555in}}{\pgfqpoint{1.694898in}{3.774256in}}%
\pgfusepath{clip}%
\pgfsetrectcap%
\pgfsetroundjoin%
\pgfsetlinewidth{1.505625pt}%
\definecolor{currentstroke}{rgb}{0.000000,0.000000,1.000000}%
\pgfsetstrokecolor{currentstroke}%
\pgfsetdash{}{0pt}%
\pgfpathmoveto{\pgfqpoint{0.704124in}{0.492112in}}%
\pgfpathlineto{\pgfqpoint{2.244940in}{2.207683in}}%
\pgfpathlineto{\pgfqpoint{1.628614in}{3.923254in}}%
\pgfusepath{stroke}%
\end{pgfscope}%
\begin{pgfscope}%
\pgfsetrectcap%
\pgfsetmiterjoin%
\pgfsetlinewidth{0.803000pt}%
\definecolor{currentstroke}{rgb}{0.000000,0.000000,0.000000}%
\pgfsetstrokecolor{currentstroke}%
\pgfsetdash{}{0pt}%
\pgfpathmoveto{\pgfqpoint{0.627083in}{0.320555in}}%
\pgfpathlineto{\pgfqpoint{0.627083in}{4.094811in}}%
\pgfusepath{stroke}%
\end{pgfscope}%
\begin{pgfscope}%
\pgfsetrectcap%
\pgfsetmiterjoin%
\pgfsetlinewidth{0.803000pt}%
\definecolor{currentstroke}{rgb}{0.000000,0.000000,0.000000}%
\pgfsetstrokecolor{currentstroke}%
\pgfsetdash{}{0pt}%
\pgfpathmoveto{\pgfqpoint{2.321981in}{0.320555in}}%
\pgfpathlineto{\pgfqpoint{2.321981in}{4.094811in}}%
\pgfusepath{stroke}%
\end{pgfscope}%
\begin{pgfscope}%
\pgfsetrectcap%
\pgfsetmiterjoin%
\pgfsetlinewidth{0.803000pt}%
\definecolor{currentstroke}{rgb}{0.000000,0.000000,0.000000}%
\pgfsetstrokecolor{currentstroke}%
\pgfsetdash{}{0pt}%
\pgfpathmoveto{\pgfqpoint{0.627083in}{0.320555in}}%
\pgfpathlineto{\pgfqpoint{2.321981in}{0.320555in}}%
\pgfusepath{stroke}%
\end{pgfscope}%
\begin{pgfscope}%
\pgfsetrectcap%
\pgfsetmiterjoin%
\pgfsetlinewidth{0.803000pt}%
\definecolor{currentstroke}{rgb}{0.000000,0.000000,0.000000}%
\pgfsetstrokecolor{currentstroke}%
\pgfsetdash{}{0pt}%
\pgfpathmoveto{\pgfqpoint{0.627083in}{4.094811in}}%
\pgfpathlineto{\pgfqpoint{2.321981in}{4.094811in}}%
\pgfusepath{stroke}%
\end{pgfscope}%
\begin{pgfscope}%
\pgfsetbuttcap%
\pgfsetmiterjoin%
\definecolor{currentfill}{rgb}{1.000000,1.000000,1.000000}%
\pgfsetfillcolor{currentfill}%
\pgfsetlinewidth{0.000000pt}%
\definecolor{currentstroke}{rgb}{0.000000,0.000000,0.000000}%
\pgfsetstrokecolor{currentstroke}%
\pgfsetstrokeopacity{0.000000}%
\pgfsetdash{}{0pt}%
\pgfpathmoveto{\pgfqpoint{2.826921in}{0.320555in}}%
\pgfpathlineto{\pgfqpoint{4.521819in}{0.320555in}}%
\pgfpathlineto{\pgfqpoint{4.521819in}{4.094811in}}%
\pgfpathlineto{\pgfqpoint{2.826921in}{4.094811in}}%
\pgfpathlineto{\pgfqpoint{2.826921in}{0.320555in}}%
\pgfpathclose%
\pgfusepath{fill}%
\end{pgfscope}%
\begin{pgfscope}%
\pgfpathrectangle{\pgfqpoint{2.826921in}{0.320555in}}{\pgfqpoint{1.694898in}{3.774256in}}%
\pgfusepath{clip}%
\pgfsetbuttcap%
\pgfsetroundjoin%
\pgfsetlinewidth{0.501875pt}%
\definecolor{currentstroke}{rgb}{0.501961,0.501961,0.501961}%
\pgfsetstrokecolor{currentstroke}%
\pgfsetstrokeopacity{0.500000}%
\pgfsetdash{{1.850000pt}{0.800000pt}}{0.000000pt}%
\pgfpathmoveto{\pgfqpoint{3.212125in}{0.320555in}}%
\pgfpathlineto{\pgfqpoint{3.212125in}{4.094811in}}%
\pgfusepath{stroke}%
\end{pgfscope}%
\begin{pgfscope}%
\pgfsetbuttcap%
\pgfsetroundjoin%
\definecolor{currentfill}{rgb}{0.000000,0.000000,0.000000}%
\pgfsetfillcolor{currentfill}%
\pgfsetlinewidth{0.803000pt}%
\definecolor{currentstroke}{rgb}{0.000000,0.000000,0.000000}%
\pgfsetstrokecolor{currentstroke}%
\pgfsetdash{}{0pt}%
\pgfsys@defobject{currentmarker}{\pgfqpoint{0.000000in}{-0.048611in}}{\pgfqpoint{0.000000in}{0.000000in}}{%
\pgfpathmoveto{\pgfqpoint{0.000000in}{0.000000in}}%
\pgfpathlineto{\pgfqpoint{0.000000in}{-0.048611in}}%
\pgfusepath{stroke,fill}%
}%
\begin{pgfscope}%
\pgfsys@transformshift{3.212125in}{0.320555in}%
\pgfsys@useobject{currentmarker}{}%
\end{pgfscope}%
\end{pgfscope}%
\begin{pgfscope}%
\definecolor{textcolor}{rgb}{0.000000,0.000000,0.000000}%
\pgfsetstrokecolor{textcolor}%
\pgfsetfillcolor{textcolor}%
\pgftext[x=3.212125in,y=0.223333in,,top]{\color{textcolor}{\rmfamily\fontsize{10.000000}{12.000000}\selectfont\catcode`\^=\active\def^{\ifmmode\sp\else\^{}\fi}\catcode`\%=\active\def%{\%}$\mathdefault{100}$}}%
\end{pgfscope}%
\begin{pgfscope}%
\pgfpathrectangle{\pgfqpoint{2.826921in}{0.320555in}}{\pgfqpoint{1.694898in}{3.774256in}}%
\pgfusepath{clip}%
\pgfsetbuttcap%
\pgfsetroundjoin%
\pgfsetlinewidth{0.501875pt}%
\definecolor{currentstroke}{rgb}{0.501961,0.501961,0.501961}%
\pgfsetstrokecolor{currentstroke}%
\pgfsetstrokeopacity{0.500000}%
\pgfsetdash{{1.850000pt}{0.800000pt}}{0.000000pt}%
\pgfpathmoveto{\pgfqpoint{3.828452in}{0.320555in}}%
\pgfpathlineto{\pgfqpoint{3.828452in}{4.094811in}}%
\pgfusepath{stroke}%
\end{pgfscope}%
\begin{pgfscope}%
\pgfsetbuttcap%
\pgfsetroundjoin%
\definecolor{currentfill}{rgb}{0.000000,0.000000,0.000000}%
\pgfsetfillcolor{currentfill}%
\pgfsetlinewidth{0.803000pt}%
\definecolor{currentstroke}{rgb}{0.000000,0.000000,0.000000}%
\pgfsetstrokecolor{currentstroke}%
\pgfsetdash{}{0pt}%
\pgfsys@defobject{currentmarker}{\pgfqpoint{0.000000in}{-0.048611in}}{\pgfqpoint{0.000000in}{0.000000in}}{%
\pgfpathmoveto{\pgfqpoint{0.000000in}{0.000000in}}%
\pgfpathlineto{\pgfqpoint{0.000000in}{-0.048611in}}%
\pgfusepath{stroke,fill}%
}%
\begin{pgfscope}%
\pgfsys@transformshift{3.828452in}{0.320555in}%
\pgfsys@useobject{currentmarker}{}%
\end{pgfscope}%
\end{pgfscope}%
\begin{pgfscope}%
\definecolor{textcolor}{rgb}{0.000000,0.000000,0.000000}%
\pgfsetstrokecolor{textcolor}%
\pgfsetfillcolor{textcolor}%
\pgftext[x=3.828452in,y=0.223333in,,top]{\color{textcolor}{\rmfamily\fontsize{10.000000}{12.000000}\selectfont\catcode`\^=\active\def^{\ifmmode\sp\else\^{}\fi}\catcode`\%=\active\def%{\%}$\mathdefault{120}$}}%
\end{pgfscope}%
\begin{pgfscope}%
\pgfpathrectangle{\pgfqpoint{2.826921in}{0.320555in}}{\pgfqpoint{1.694898in}{3.774256in}}%
\pgfusepath{clip}%
\pgfsetbuttcap%
\pgfsetroundjoin%
\pgfsetlinewidth{0.501875pt}%
\definecolor{currentstroke}{rgb}{0.501961,0.501961,0.501961}%
\pgfsetstrokecolor{currentstroke}%
\pgfsetstrokeopacity{0.500000}%
\pgfsetdash{{1.850000pt}{0.800000pt}}{0.000000pt}%
\pgfpathmoveto{\pgfqpoint{4.444778in}{0.320555in}}%
\pgfpathlineto{\pgfqpoint{4.444778in}{4.094811in}}%
\pgfusepath{stroke}%
\end{pgfscope}%
\begin{pgfscope}%
\pgfsetbuttcap%
\pgfsetroundjoin%
\definecolor{currentfill}{rgb}{0.000000,0.000000,0.000000}%
\pgfsetfillcolor{currentfill}%
\pgfsetlinewidth{0.803000pt}%
\definecolor{currentstroke}{rgb}{0.000000,0.000000,0.000000}%
\pgfsetstrokecolor{currentstroke}%
\pgfsetdash{}{0pt}%
\pgfsys@defobject{currentmarker}{\pgfqpoint{0.000000in}{-0.048611in}}{\pgfqpoint{0.000000in}{0.000000in}}{%
\pgfpathmoveto{\pgfqpoint{0.000000in}{0.000000in}}%
\pgfpathlineto{\pgfqpoint{0.000000in}{-0.048611in}}%
\pgfusepath{stroke,fill}%
}%
\begin{pgfscope}%
\pgfsys@transformshift{4.444778in}{0.320555in}%
\pgfsys@useobject{currentmarker}{}%
\end{pgfscope}%
\end{pgfscope}%
\begin{pgfscope}%
\definecolor{textcolor}{rgb}{0.000000,0.000000,0.000000}%
\pgfsetstrokecolor{textcolor}%
\pgfsetfillcolor{textcolor}%
\pgftext[x=4.444778in,y=0.223333in,,top]{\color{textcolor}{\rmfamily\fontsize{10.000000}{12.000000}\selectfont\catcode`\^=\active\def^{\ifmmode\sp\else\^{}\fi}\catcode`\%=\active\def%{\%}$\mathdefault{140}$}}%
\end{pgfscope}%
\begin{pgfscope}%
\pgfpathrectangle{\pgfqpoint{2.826921in}{0.320555in}}{\pgfqpoint{1.694898in}{3.774256in}}%
\pgfusepath{clip}%
\pgfsetbuttcap%
\pgfsetroundjoin%
\pgfsetlinewidth{0.501875pt}%
\definecolor{currentstroke}{rgb}{0.501961,0.501961,0.501961}%
\pgfsetstrokecolor{currentstroke}%
\pgfsetstrokeopacity{0.500000}%
\pgfsetdash{{1.850000pt}{0.800000pt}}{0.000000pt}%
\pgfpathmoveto{\pgfqpoint{2.826921in}{0.492112in}}%
\pgfpathlineto{\pgfqpoint{4.521819in}{0.492112in}}%
\pgfusepath{stroke}%
\end{pgfscope}%
\begin{pgfscope}%
\pgfsetbuttcap%
\pgfsetroundjoin%
\definecolor{currentfill}{rgb}{0.000000,0.000000,0.000000}%
\pgfsetfillcolor{currentfill}%
\pgfsetlinewidth{0.803000pt}%
\definecolor{currentstroke}{rgb}{0.000000,0.000000,0.000000}%
\pgfsetstrokecolor{currentstroke}%
\pgfsetdash{}{0pt}%
\pgfsys@defobject{currentmarker}{\pgfqpoint{-0.048611in}{0.000000in}}{\pgfqpoint{-0.000000in}{0.000000in}}{%
\pgfpathmoveto{\pgfqpoint{-0.000000in}{0.000000in}}%
\pgfpathlineto{\pgfqpoint{-0.048611in}{0.000000in}}%
\pgfusepath{stroke,fill}%
}%
\begin{pgfscope}%
\pgfsys@transformshift{2.826921in}{0.492112in}%
\pgfsys@useobject{currentmarker}{}%
\end{pgfscope}%
\end{pgfscope}%
\begin{pgfscope}%
\pgfpathrectangle{\pgfqpoint{2.826921in}{0.320555in}}{\pgfqpoint{1.694898in}{3.774256in}}%
\pgfusepath{clip}%
\pgfsetbuttcap%
\pgfsetroundjoin%
\pgfsetlinewidth{0.501875pt}%
\definecolor{currentstroke}{rgb}{0.501961,0.501961,0.501961}%
\pgfsetstrokecolor{currentstroke}%
\pgfsetstrokeopacity{0.500000}%
\pgfsetdash{{1.850000pt}{0.800000pt}}{0.000000pt}%
\pgfpathmoveto{\pgfqpoint{2.826921in}{2.207683in}}%
\pgfpathlineto{\pgfqpoint{4.521819in}{2.207683in}}%
\pgfusepath{stroke}%
\end{pgfscope}%
\begin{pgfscope}%
\pgfsetbuttcap%
\pgfsetroundjoin%
\definecolor{currentfill}{rgb}{0.000000,0.000000,0.000000}%
\pgfsetfillcolor{currentfill}%
\pgfsetlinewidth{0.803000pt}%
\definecolor{currentstroke}{rgb}{0.000000,0.000000,0.000000}%
\pgfsetstrokecolor{currentstroke}%
\pgfsetdash{}{0pt}%
\pgfsys@defobject{currentmarker}{\pgfqpoint{-0.048611in}{0.000000in}}{\pgfqpoint{-0.000000in}{0.000000in}}{%
\pgfpathmoveto{\pgfqpoint{-0.000000in}{0.000000in}}%
\pgfpathlineto{\pgfqpoint{-0.048611in}{0.000000in}}%
\pgfusepath{stroke,fill}%
}%
\begin{pgfscope}%
\pgfsys@transformshift{2.826921in}{2.207683in}%
\pgfsys@useobject{currentmarker}{}%
\end{pgfscope}%
\end{pgfscope}%
\begin{pgfscope}%
\pgfpathrectangle{\pgfqpoint{2.826921in}{0.320555in}}{\pgfqpoint{1.694898in}{3.774256in}}%
\pgfusepath{clip}%
\pgfsetbuttcap%
\pgfsetroundjoin%
\pgfsetlinewidth{0.501875pt}%
\definecolor{currentstroke}{rgb}{0.501961,0.501961,0.501961}%
\pgfsetstrokecolor{currentstroke}%
\pgfsetstrokeopacity{0.500000}%
\pgfsetdash{{1.850000pt}{0.800000pt}}{0.000000pt}%
\pgfpathmoveto{\pgfqpoint{2.826921in}{3.923254in}}%
\pgfpathlineto{\pgfqpoint{4.521819in}{3.923254in}}%
\pgfusepath{stroke}%
\end{pgfscope}%
\begin{pgfscope}%
\pgfsetbuttcap%
\pgfsetroundjoin%
\definecolor{currentfill}{rgb}{0.000000,0.000000,0.000000}%
\pgfsetfillcolor{currentfill}%
\pgfsetlinewidth{0.803000pt}%
\definecolor{currentstroke}{rgb}{0.000000,0.000000,0.000000}%
\pgfsetstrokecolor{currentstroke}%
\pgfsetdash{}{0pt}%
\pgfsys@defobject{currentmarker}{\pgfqpoint{-0.048611in}{0.000000in}}{\pgfqpoint{-0.000000in}{0.000000in}}{%
\pgfpathmoveto{\pgfqpoint{-0.000000in}{0.000000in}}%
\pgfpathlineto{\pgfqpoint{-0.048611in}{0.000000in}}%
\pgfusepath{stroke,fill}%
}%
\begin{pgfscope}%
\pgfsys@transformshift{2.826921in}{3.923254in}%
\pgfsys@useobject{currentmarker}{}%
\end{pgfscope}%
\end{pgfscope}%
\begin{pgfscope}%
\definecolor{textcolor}{rgb}{0.000000,0.000000,0.000000}%
\pgfsetstrokecolor{textcolor}%
\pgfsetfillcolor{textcolor}%
\pgftext[x=2.722754in,y=2.207683in,,bottom,rotate=90.000000]{\color{textcolor}{\rmfamily\fontsize{10.000000}{12.000000}\selectfont\catcode`\^=\active\def^{\ifmmode\sp\else\^{}\fi}\catcode`\%=\active\def%{\%}Producto B}}%
\end{pgfscope}%
\begin{pgfscope}%
\pgfpathrectangle{\pgfqpoint{2.826921in}{0.320555in}}{\pgfqpoint{1.694898in}{3.774256in}}%
\pgfusepath{clip}%
\pgfsetrectcap%
\pgfsetroundjoin%
\pgfsetlinewidth{1.505625pt}%
\definecolor{currentstroke}{rgb}{0.000000,0.501961,0.000000}%
\pgfsetstrokecolor{currentstroke}%
\pgfsetdash{}{0pt}%
\pgfpathmoveto{\pgfqpoint{2.903962in}{0.492112in}}%
\pgfpathlineto{\pgfqpoint{3.520288in}{2.207683in}}%
\pgfpathlineto{\pgfqpoint{4.444778in}{3.923254in}}%
\pgfusepath{stroke}%
\end{pgfscope}%
\begin{pgfscope}%
\pgfsetrectcap%
\pgfsetmiterjoin%
\pgfsetlinewidth{0.803000pt}%
\definecolor{currentstroke}{rgb}{0.000000,0.000000,0.000000}%
\pgfsetstrokecolor{currentstroke}%
\pgfsetdash{}{0pt}%
\pgfpathmoveto{\pgfqpoint{2.826921in}{0.320555in}}%
\pgfpathlineto{\pgfqpoint{2.826921in}{4.094811in}}%
\pgfusepath{stroke}%
\end{pgfscope}%
\begin{pgfscope}%
\pgfsetrectcap%
\pgfsetmiterjoin%
\pgfsetlinewidth{0.803000pt}%
\definecolor{currentstroke}{rgb}{0.000000,0.000000,0.000000}%
\pgfsetstrokecolor{currentstroke}%
\pgfsetdash{}{0pt}%
\pgfpathmoveto{\pgfqpoint{4.521819in}{0.320555in}}%
\pgfpathlineto{\pgfqpoint{4.521819in}{4.094811in}}%
\pgfusepath{stroke}%
\end{pgfscope}%
\begin{pgfscope}%
\pgfsetrectcap%
\pgfsetmiterjoin%
\pgfsetlinewidth{0.803000pt}%
\definecolor{currentstroke}{rgb}{0.000000,0.000000,0.000000}%
\pgfsetstrokecolor{currentstroke}%
\pgfsetdash{}{0pt}%
\pgfpathmoveto{\pgfqpoint{2.826921in}{0.320555in}}%
\pgfpathlineto{\pgfqpoint{4.521819in}{0.320555in}}%
\pgfusepath{stroke}%
\end{pgfscope}%
\begin{pgfscope}%
\pgfsetrectcap%
\pgfsetmiterjoin%
\pgfsetlinewidth{0.803000pt}%
\definecolor{currentstroke}{rgb}{0.000000,0.000000,0.000000}%
\pgfsetstrokecolor{currentstroke}%
\pgfsetdash{}{0pt}%
\pgfpathmoveto{\pgfqpoint{2.826921in}{4.094811in}}%
\pgfpathlineto{\pgfqpoint{4.521819in}{4.094811in}}%
\pgfusepath{stroke}%
\end{pgfscope}%
\begin{pgfscope}%
\definecolor{textcolor}{rgb}{0.000000,0.000000,0.000000}%
\pgfsetstrokecolor{textcolor}%
\pgfsetfillcolor{textcolor}%
\pgftext[x=2.323312in,y=4.722560in,,top]{\color{textcolor}{\rmfamily\fontsize{14.000000}{16.800000}\selectfont\catcode`\^=\active\def^{\ifmmode\sp\else\^{}\fi}\catcode`\%=\active\def%{\%}Gráfico de ejes compartidos con personalización}}%
\end{pgfscope}%
\end{pgfpicture}%
\makeatother%
\endgroup%
 
    \caption{Gráfico Personalizado con rcParams}
    \label{fig:grafico_objeto_con_rcParams}
\end{figure}

\begin{nota}
    Si observo detenidamente se utiliza \texttt{axes.grid} ya que \textit{Matplot} organiza sus parámetros por categorías. La grilla por defecto pertenece a la categoría \texttt{axes}.

    ¿Por qué utilizo \texttt{rcParams.update()}? Si utilizara \texttt{rcParams['parámetro'] = 'estilo'} deberia utilizar esa misma linea para modificar cada parámetro. En cambio si agregamos \texttt{.update} con un diccionario, podemos modificar varios parámetros de una sola vez.
\end{nota}

\begin{nota}
    En mi configuración para \LaTeX  \texttt{"text.usetex": True}, le pido a \textit{Python} que llame a la instalación local de \LaTeX (\textit{TeX Live} o \textit{MiKTeX}) para renderizar cada letra del gráfico. Esto garantiza que fórmulas matemáticas y fuentes (\texttt{CMU Serif}) coincidan perfectamente con tu documento principal.
\end{nota}

\begin{nota}
    \textit{Orden de precedencia:} Si usted define algo en \texttt{rcParams} pero luego, define un parametro dentro de un \texttt{ax.plot(color='red')}, que en este caso especifica un color manual, el comando local (\texttt{ax.plot()}) gana sobre la configuración global.
\end{nota}



\subsection{Creación de temas personalizados (\texttt{.mplstyle})}

Para crear una modificación personalizada que pueda aplicar a cualquier proyecto en el cual utilice \textit{Matplotlib} debe crear un archivo \texttt{.mplstyle}. Este documento puede tener cualquier nombre, pero terminar con el formato .\texttt{mplstyle}.

Dentro del archivo escriba las configuraciones utilizando la lógica de \texttt{key:Value}, pero sin comillas y usando ``\#'' para los comentarios.

por ej:
\begin{pycode}{Archivo .mplstyle}
# Estilo Personalizado para Reportes

# FUENTES Y TEXTOS
font.family : serif
font.serif  : DejaVu Serif, Times New Roman
font.size   : 11
axes.titlesize : 14
axes.titleweight : bold
axes.labelsize : 12

# LINEAS Y MARCADORES
lines.linewidth : 2.5
lines.markersize : 8
lines.marker : o

# OTROS...
\end{pycode}

\subsubsection{Ubicación y Carga del Tema}

Para que \textit{Matplotlib} reconozca su estilo de forma global (desde cualquier proyecto), debe guardarlo en la carpeta de configuración del sistema:

\begin{itemize}
\item \textit{Windows:} \texttt{C:\textbackslash Usuarios\textbackslash Usuario\textbackslash .matplotlib\textbackslash stylelib\textbackslash}
\item \textit{Linux/Mac:} \texttt{\textasciitilde/.config/matplotlib/stylelib/}
\end{itemize}

Si lo guarda allí, podrá invocarlo simplemente por su nombre:
\begin{verbatim}
plt.style.use('mi_estilo')
\end{verbatim}

Si prefiere mantener el estilo dentro de la carpeta de un proyecto específico, utilice la ruta relativa:
\begin{verbatim}
plt.style.use('./estilo_proyecto.mplstyle')
\end{verbatim}




\subsection{Gestión de Ciclos de Color (\texttt{prop\_cycle})}

Cuando graficamos múltiples series de datos en un mismo objeto \texttt{ax}, \textit{Matplotlib} debe decidir qué color asignarle a cada una para que no se confundan. Esta decisión la toma consultando el \texttt{prop\_cycle} (ciclo de propiedades).

\subsubsection{Personalización con \texttt{cycler}}

Para crear una experiencia visual profesional, podemos definir no solo colores, sino también combinar estilos de línea y marcadores en un solo ciclo. Esto asegura que, incluso si el reporte se imprime en blanco y negro, las líneas sean distinguibles por su trazo (sólido, punteado, etc.).

\begin{pycode}{Configuración de ciclo avanzado}
import matplotlib.pyplot as plt
from cycler import cycler

colores = ['#1f77b4', '#ff7f0e', '#2ca02c'] # Azul, Naranja, Verde
estilos = ['-', '--', ':'] # Sólido, Rayado, Punteado

# Creamos el ciclo combinado 
mi_ciclo = cycler(color=colores) + cycler(linestyle=estilos)

# Aplicamos el ciclo a la categoría 'axes'
plt.rc('axes', prop_cycle=mi_ciclo)
\end{pycode}

\subsubsection{Integración en archivos \texttt{.mplstyle}}

Para que su paleta de colores sea permanente en todos sus proyectos, puede definirla directamente en su archivo de estilo. La sintaxis dentro de un \texttt{.mplstyle} requiere que la función \texttt{cycler} se escriba en una sola línea:

\begin{verbatim}
En su archivo .mplstyle

axes.prop_cycle : cycler(
    'color', ['2980B9', '27AE60', 'E67E22', '8E44AD']
    )
\end{verbatim}
